\begin{recipe}
[% 
    preparationtime = {\unit[50]{min}},
    portion = {\portion{2}},
    source = {\protect\recipesource{https://www.zuckerjagdwurst.com/de/rezepte/vegane-ente-mit-orangensosse-und-mandelreis}{Zucker\&Jagdwurst: Vegane Ente mit Orangensoße und Mandelreis}},
    calory = {1300kcal | 52g Protein | 143g KH | 62g Fett},
    price = {6,80€},
    created = {10.06.2026},
    %rating = {1},
]
{Vegane Ente mit Orangensoße}

\graph{% pictures
    big=pic/hauptgericht/17_vegane-ente-orangensosse
}

\introduction{%
Vegane Ente aus Seitan mit knusprigem Reispapier-Mantel, dazu eine fruchtige Orangensoße und Mandelreis. Klingt nach Feiertagsessen, ist aber erstaunlich machbar.
}

\ingredients{
    \textbf{Seitan} & \\
    240 g & Seitan am Stück\index{Tofu \& Fleischalternativen!Seitan}\\
    500 ml & Heiße Gemüsebrühe\\
    2 EL & Sojasoße\\
    \textbf{Marinade und Haut} & \\
    4 große Blätter & Reispapier\index{Getreide!Reispapier}\\
    6 EL & Asiatische Pilz-Würzsoße oder vegane Austernsoße\\
    2 EL & Sojasoße\\
    2 EL & Olivenöl\\
    2 TL & Senf\\
    1 TL & Getrockneter Majoran\\
    1 TL & Getrockneter Thymian\\
    \nicefrac{1}{4} TL & Knoblauchpulver\\
    1 TL & Zwiebelpulver\\
    & Salz, Pfeffer\\
    & Pflanzenöl zum Braten\\
    \textbf{Orangensoße} & \\
    2 & Schalotten\index{Gemüse!Schalotte}\\
    6 & Orangen\index{Obst!Orange}\\
    2 EL & Zucker\\
    1 EL & Wasser\\
    2 EL & Rotweinessig\\
    250 ml & Gemüsebrühe\\
    50 g & Vegane Butter\\
    & Salz, Pfeffer\\
    & Speisestärke nach Bedarf\\
    \textbf{Mandelreis} & \\
    150 g & Basmatireis\index{Getreide!Basmatireis}\\
    75 g & Gehobelte Mandeln\index{Nüsse \& Samen!Mandel}
}

\preparation{%
    \step%
    Gemüsebrühe und Sojasoße in einem kleinen Topf aufkochen. Seitan in dicke Scheiben schneiden und etwa 10 Minuten darin köcheln lassen. Danach herausnehmen und abkühlen lassen.
    
    \step%
    Für die Marinade Pilz-Würzsoße, Sojasoße, Olivenöl, Senf und Gewürze verrühren. Mit Salz und Pfeffer abschmecken.
    
    \step%
    Für die Orangensoße Schalotten fein würfeln. Von einer Orange die Schale abreiben. Zwei Orangen filetieren und beiseitestellen. Die restlichen Orangen auspressen, bis etwa 200 ml Orangensaft entstehen.
    
    \step%
    Zucker, Wasser und Rotweinessig in einem Topf erhitzen, bis der Zucker geschmolzen ist. Schalotten zugeben und kurz glasig anschwitzen.
    
    \step%
    Orangensaft und Gemüsebrühe zugeben, aufkochen und bei niedriger Hitze köcheln lassen, während der Rest vorbereitet wird.
    
    \step%
    Reis nach Packungsanweisung kochen. Mandeln in einer Pfanne ohne Fett goldbraun rösten und anschließend unter den fertigen Reis mischen.
    
    \step%
    Die abgekühlten Seitanscheiben in der Marinade wenden. Restliche Marinade aufheben.
    
    \step%
    Reispapierblätter kurz in Wasser einweichen. Jeweils eine marinierte Seitanscheibe mittig darauflegen und einwickeln.
    
    \step%
    Seitanscheiben in einer Pfanne mit Pflanzenöl bei mittlerer Hitze von beiden Seiten anbraten, bis das Reispapier leicht knusprig und gebräunt ist. Nach Bedarf mit etwas restlicher Marinade bestreichen.
    
    \step%
    Orangensoße vom Herd nehmen und vegane Butter einrühren. Orangenfilets zugeben und mit Salz und Pfeffer abschmecken.
    
    \step%
    Falls die Soße zu dünn ist, etwas Speisestärke mit Wasser glatt rühren, einrühren und kurz köcheln lassen.
    
    \step%
    Seitan mit Orangensoße und Mandelreis servieren.
}

\hint{
    Reispapier nicht zu lange einweichen und beim Braten nicht zu heiß werden lassen. Sonst klebt es erst wie Tapete und verbrennt dann beleidigt.
}

\suggestion[Timing]{%
    Erst Soße ansetzen, dann Reis kochen und Mandeln rösten. Den Seitan am besten zuletzt braten, damit die Reispapier-Haut beim Servieren noch knusprig ist.
}

\end{recipe}